\documentclass[11pt]{article}
\usepackage[margin=1in]{geometry}
\usepackage{booktabs}
\usepackage{amsmath}
\usepackage{url}
\newcommand{\coderepo}{\url{https://github.com/jakehemmerle/mapa_diver_ensemble}}

\title{MAPA+DIVER-1: A Pre-specified Rank-Normalized Ensemble of Frozen Encoders\\for Neuroprobe Within-Session Decoding}
\author{Jake Hemmerle\\Spark}
\date{November 2026}

\begin{document}
\maketitle

\begin{abstract}
We describe the ``MAPA+DIVER-1 ensemble'' submission to the Neuroprobe leaderboard (within-session regime). Two publicly released, frozen intracranial encoders---MAPA (last layer, closed-form ridge readout) and DIVER-1 (0.1\,s tiny, linear head)---are combined by an equal-weight mean of per-member scores after each member's scores are mapped through the empirical CDF of its own label-free validation scores. No pretraining was performed. The design was fixed using only non-evaluation sessions, vetoed once on half of the official validation split, and then scored once on the official test split. On the official test split the ensemble reaches mean AUROC $0.7037$ versus $0.6955$ for our MAPA reproduction ($\Delta=+0.0082$, 6/6 subjects improved, exact sign-flip $p=0.0156$, session bootstrap 95\% CI $[+0.0036,+0.0132]$). Code: \coderepo.
\end{abstract}

\section{Introduction}
Neuroprobe~\cite{neuroprobe,neuroprobe-site,neuroprobe-repo} evaluates decoding of 15 audio, visual and language features from intracranial recordings of the Brain Treebank~\cite{btbank}. Recent pretrained encoders, MAPA~\cite{mapa} and DIVER-1~\cite{diver}, report within-session leaderboard scores of $0.6953$ and $0.6777$ (DIVER-1 frozen), respectively. The two models were pretrained on different data and use different architectures, which suggests their errors are partly decorrelated. This note documents a minimal, pre-specified combination of the two frozen encoders, the three-tier protocol used to decide whether to submit it, and how to reproduce every reported number. Its contribution is the protocol and evidence, not a new model.

\section{Method}
\subsection{Members}
\paragraph{MAPA.} We use the released MAPA encoder frozen, taking last-layer features (\texttt{mapa12}). Preprocessing reuses MAPA's released frontend unchanged (\path{ensemble/mapa/prepare_session.py}, \path{encode_session.py}): within-shaft mean re-referencing, whole-session zero-phase 0.5\,Hz FIR high-pass and 60--300\,Hz notch filtering, centered STFTs over the continuous session (the first Slow 500\,ms frame is centered on onset and thus reaches 250\,ms pre-onset), whole-session per-contact/per-frequency median/MAD normalization (computed without labels across the entire recording, including the test half), and MAPA's frozen Guard~1 per-session contact exclusions (removing 25 of 1,414 evaluation contacts). The readout is a closed-form, trace-normalized ridge regression (\path{common/ridge.py}) on features standardized with \emph{training rows only} (\path{ensemble/mapa/gram.py}); a single penalty per cell is chosen by validation AUROC from a 25-point log-spaced grid of multipliers.
\paragraph{DIVER-1.} We use the released DIVER-1 0.1\,s \emph{tiny} encoder frozen, with features in bfloat16. Preprocessing matches DIVER-1's upstream \path{preprocessing_for_neuroprobe.py}: whole-session zero-phase 0.5\,Hz high-pass and 60/120/180\,Hz notch filtering and $\pm 200\,\mu\text{V}$ clipping, followed by 0--1\,s window slicing, Laplacian re-referencing, and downsampling to 500\,Hz. A linear head is trained with DIVER's own \texttt{Trainer}: AdamW (lr $0.002$, weight decay $0.01$), 40-epoch cosine schedule, early stopping with patience 10 on validation AUROC (\path{ensemble/members/diver.py}). On Tier~P we ran three seeds (0--2) to measure seed spread; Tier~V and Tier~T confirmation ran seed~0 only, omitting DIVER's five out-of-fold training refits (which uniform rung~(a) never uses) after verifying exact $0.0$ score and AUROC parity on Tier~P (\path{results/parity/diver_uniform_confirmation.json}).
\paragraph{Ridge6 (reference only).} A classical spectral-feature ridge model (reimplemented from Neuroprobe PR~\#20) was computed alongside as a reference member; it is not part of the submitted ensemble.

\subsection{Fusion}
Let $s_m(x)$ be member $m$'s score and $V_m$ its scores on the label-free validation inputs of the same cell. The ensemble score is
\begin{equation}
  s(x) = \frac{1}{2}\sum_{m\in\{\text{MAPA},\text{DIVER}\}} \hat F_{V_m}\big(s_m(x)\big),
\end{equation}
where $\hat F_{V_m}$ is the empirical CDF of $V_m$ (\texttt{ensemble/fuse.py}, \texttt{NORMALIZATION = "select\_cdf"}). The weights are fixed (equal); no stacking or learned weighting is used. Validation \emph{labels} are not used by the fusion.

\section{Evaluation protocol}
The evaluation tiers, thresholds, controls, and confirmation lock below were specified in dated commits to our internal repository before the corresponding data were scored; they were not lodged with a public registry. Two post-hoc adjustments occurred during Tier~P design before the candidate was frozen and before any evaluation-session label was opened (disclosed in Sec.~\ref{sec:disclosures}); the Tier~V veto and single Tier~T scoring ran unchanged against the committed freeze.
\paragraph{Tiers.} \emph{Tier P}: 13 permitted non-evaluation Brain Treebank sessions (\texttt{1:0}, \texttt{2:1}, \texttt{2:2}, \texttt{2:3}, \texttt{2:5}, \texttt{2:6}, \texttt{3:2}, \texttt{4:2}, \texttt{5:0}, \texttt{6:1}, \texttt{6:4}, \texttt{8:0}, \texttt{9:0}), used for every design decision (member choice, layer, fusion rule, grids). Permitted session \texttt{6:0} (\emph{Megamind}, shown in evaluation sessions \texttt{4:1} and \texttt{7:1}) and partial sessions \texttt{7:100--102}, \texttt{10:100--101} (evaluation films) were excluded in code (\texttt{ensemble/paths.py}). \emph{Tier V}: the 12 evaluation sessions; penalty selection and early stopping use the first half of the official validation split, and the frozen design is scored once on the second half (2\,s purge between halves) as a veto. \emph{Tier T}: train on official train, select on official validation, score once on official test; test labels were locked until Tier V returned Go.
\paragraph{Unit and statistics.} The unit of analysis is the subject. We report the subject-weighted mean AUROC, the number of subjects improved, an exact one-sided sign-flip test over subjects, a session-level paired bootstrap 95\% CI, and leave-one-subject-out (LOSO) deltas.
\paragraph{Thresholds (pre-specified).} Tier P: Go if $\Delta\ge+0.010$ and $\ge 7/8$ subjects improve; Hold if $+0.005\le\Delta<+0.010$. Tier V: Go if $\Delta\ge+0.005$ and $\ge5/6$ subjects improve, otherwise veto. Tier T is reported regardless of outcome.
\paragraph{Controls.} Classical-feature candidates were compared with a random-projection control and a sham control of matched dimension. The ensemble candidate instead ran the ensemble controls of \texttt{ensemble/controls.py} (plan phase E3) on Tier P, all of which it passed: (i)~\emph{self-ensemble}: MAPA refit on five bootstrap resamples of its training trials and averaged scores $0.68976$, a change of $-0.00618$ versus MAPA, so the ensemble's gain exceeds what variance reduction alone buys by $0.01164$ (required margin $0.003$); (ii)~\emph{leave-one-member-out}: removing DIVER-1 costs $0.00546$ and removing MAPA costs $0.02399$, both above the $0.001$ drop threshold, so no member was dropped; (iii)~\emph{shuffled stacker}: recorded as passed (tolerance $0.002$; the frozen design uses the fixed-weight rung (a), not a stacker); (iv)~\emph{subject concentration}: the gain is carried by no single subject, with leave-one-subject-out deltas from $+0.0037$ (without subject~5) to $+0.0073$.

\section{Results}
\subsection{Tier P (design)}
\begin{table}[h]\centering\small
\begin{tabular}{lcccccl}\toprule
Candidate & Score & MAPA & $\Delta$ & Subj.\ up & $p$ & Decision\\\midrule
MAPA+DIVER-1 (submitted) & 0.70140 & 0.69594 & $+0.00546$ & 6/8 & 0.059 & Hold\\
MAPA+DIVER-1+Ridge6 & 0.70118 & 0.69594 & $+0.00524$ & 6/8 & 0.082 & Hold\\
MAPA+Ridge6 & 0.69585 & 0.69594 & $-0.00009$ & 5/8 & 0.512 & Stop\\
MAPA + classical features & -- & -- & $+0.000039$ & -- & -- & Stop\\\bottomrule
\end{tabular}
\caption{Tier P within-session (13 sessions, 8 subjects, 390 cells). Bootstrap CI for the submitted ensemble: $[-0.0013,+0.0097]$. Frozen as config \texttt{ensemble:cc49ff955fa0ae88}.}
\end{table}
Because MAPA was pretrained on 12 of the 13 Tier~P sessions (all except \texttt{5:0}), we report \texttt{5:0} separately as planned: on held-out session \texttt{5:0} the ensemble gained $+0.0180$ ($0.7306$ vs.\ MAPA $0.7126$), compared with $+0.0037$ across the 12 MAPA-pretrained Tier~P sessions (7 subjects). Across three DIVER-1 seeds on Tier~P, DIVER-1 macro AUROC was $0.6774$, $0.6785$, and $0.6789$ ($\text{SD}=0.00063$). In phase~E2, the out-of-fold logistic stacker (rung~(b)) added $+0.00226$ over uniform rung~(a) (below the $+0.003$ adoption threshold; shuffled-label stacker delta $-0.00003$), so fixed-weight rung~(a) was retained.

\subsection{Amendment before held-out scoring}\label{sec:amendment}
The first Tier V run blocked: in session 2:0, task \emph{volume}, fold~1, all 437 selection-half validation rows had label 0, so no penalty or early-stopping criterion was defined. Before any held-out score was computed, we pre-specified a fallback derived only from Tier P medians: MAPA ridge multiplier $0.46415888336127775$, Ridge6 multiplier $2.154434690031882$, DIVER 8 completed epochs (the 40-epoch cosine schedule kept and truncated). On practice data the amended code reproduced ordinary-case predictions exactly (12 checks). The blocked receipt is preserved and the amended run writes to its own namespace (protocol \texttt{within-single-class-select-v1}). On Tier~T, $0/360$ cells triggered the fallback; it applied only to nine cells during Tier~V.

\subsection{Tier V (veto)}
Ensemble $0.71110$ vs.\ MAPA $0.70277$: $\Delta=+0.00834$, 5/6 subjects improved, $p=0.047$, CI $[+0.0049,+0.0124]$, \textbf{Go}. Nine cells (volume or frame\_brightness, fold~1) had a single-class evaluation half and are undefined; 351 cells were scored. Subject~1 was negative ($-0.0049$); subject~3 was flat ($+0.0001$).

\subsection{Tier T (final, run once)}
\begin{table}[h]\centering\small
\begin{tabular}{lc}\toprule
Model & Mean test AUROC\\\midrule
MAPA+DIVER-1 ensemble & \textbf{0.70374}\\
MAPA (our reproduction) & 0.69554\\
DIVER-1 (our member) & 0.67907\\
Ridge6 (reference) & 0.67602\\\midrule
MAPA (leaderboard) & 0.6953\\
DIVER-1 frozen (leaderboard) & 0.6777\\\bottomrule
\end{tabular}
\caption{Tier T within-session, 12 sessions, 360/360 cells, no undefined cells. Ensemble vs.\ MAPA: $\Delta=+0.00820$, 6/6 subjects, $p=0.0156$, CI $[+0.0036,+0.0132]$, minimum LOSO $\Delta=+0.0064$.}
\end{table}
The leaderboard metric (mean over tasks of the mean over all session$\times$fold test AUROCs) equals our subject-weighted score here because the Tier T grid is balanced and complete.

\begin{table}[h]\centering\small
\begin{tabular}{lr@{\qquad}lr@{\qquad}lr}\toprule
Session & $\Delta$ & Session & $\Delta$ & Session & $\Delta$\\\midrule
1:1 & $+0.0012$ & 3:0 & $+0.0067$ & 7:0 & $+0.0044$\\
1:2 & $+0.0042$ & 3:1 & $-0.0044$ & 7:1 & $+0.0297$\\
2:0 & $+0.0144$ & 4:0 & $+0.0176$ & 10:0 & $+0.0112$\\
2:4 & $+0.0081$ & 4:1 & $-0.0021$ & 10:1 & $+0.0074$\\\bottomrule
\end{tabular}
\caption{Tier T per-session deltas (ensemble $-$ MAPA). Per subject: 1: $+0.0027$, 2: $+0.0112$, 3: $+0.0011$, 4: $+0.0078$, 7: $+0.0170$, 10: $+0.0093$.}
\end{table}

\begin{table}[h]\centering\small
\begin{tabular}{lr@{\qquad}lr}\toprule
Task & $\Delta$ & Task & $\Delta$\\\midrule
word\_index & $+0.0345$ & speech & $+0.0057$\\
delta\_volume & $+0.0216$ & pitch & $+0.0008$\\
word\_length & $+0.0206$ & global\_flow & $-0.0022$\\
word\_part\_speech & $+0.0180$ & frame\_brightness & $-0.0040$\\
word\_gap & $+0.0169$ & face\_num & $-0.0064$\\
gpt2\_surprisal & $+0.0139$ & local\_flow & $-0.0090$\\
onset & $+0.0131$ & volume & $-0.0106$\\
word\_head\_pos & $+0.0103$ & & \\\bottomrule
\end{tabular}
\caption{Tier T per-task deltas. No task is significant after Holm correction (smallest $p_{\mathrm{Holm}}=0.234$).}
\end{table}

\section{Reproduction}
Code: \coderepo. The public repository is a snapshot of the code that produced Tier T (internal commit \texttt{50932eb}); Tier V ran at internal commit \texttt{95b94ba}, which differs only in how many DIVER-1 export workers run per GPU. Both runs use frozen config \texttt{cc49ff955fa0ae88}.
\begin{enumerate}
\item Install with \texttt{uv sync} (\path{pyproject.toml}, \path{uv.lock}, recorded from the server environment used for every run); obtain the Brain Treebank, Neuroprobe~\cite{neuroprobe-repo}, and the released MAPA and DIVER-1 0.1\,s tiny checkpoints. Our \path{common/splits.py} calls \path{neuroprobe.train_test_splits.generate_splits_within_session} directly at revision \texttt{f9b0842} and cross-checks every fold against its exact split arithmetic; current Neuroprobe \texttt{main} is \texttt{b282acd}, which differs from \texttt{f9b0842} only in leaderboard files.
\item Tier P design: \path{ensemble/} and \path{classical/} experiment drivers; freeze with \path{server/overnight_freeze.py}, which writes \path{results/frozen/within.json} and attests its hash in Git.
\item Confirmation: \texttt{python -m server.overnight\_confirm within} runs Tier V (veto) and, only on Go, Tier T against the lock file. The single-class fallback is \path{common/confirm_fallback.py} and \path{server/confirm_amendment.py}.
\item Components: \path{common/ridge.py} (ridge readout), \path{ensemble/members/diver.py} (DIVER head), \path{ensemble/fuse.py} (fusion), \path{common/splits.py} (tier splits, purge), \path{common/stats.py} (sign-flip, bootstrap, LOSO).
\item Audit artifacts (committed in the public repository under \path{results/}): \path{results/ledger.csv} (hash-chained, append-only, 381 rows), \path{results/test_access.log}, \path{results/overnight/confirmation_V_within.json}, \path{results/overnight/confirmation_T_within.json}, \path{results/overnight/confirmation_blocked_within.json}, \path{results/overnight/amendment_within_single_class_v1.json}, \path{results/ensemble/e1_within_w2.json} through \path{e5_within.json}, \path{results/parity/diver_uniform_confirmation.json}, and \path{results/submission/export_t_within.py}.
\end{enumerate}
Compute: one server with 4$\times$ NVIDIA CMP~170HX. The Tier T DIVER-1 feature export ran as 6 workers from 17:58 to 19:41 CST on 2026-11-26 (about 1\,h\,43\,min wall clock, per \path{logs/e0-w2-T.within-diver-*.log}).

\section{Disclosures and limitations}\label{sec:disclosures}
\begin{itemize}
\item \textbf{Tier P protocol adjustments (before freeze).} Two adjustments were made on Tier~P before freezing the candidate and before any evaluation-session label was opened (recorded in \path{06-classical-and-ensemble-plan.md}):
  (1)~\emph{Fusion normalization:} the initial Wave-1 screen (\texttt{mapa12}, \texttt{mapa0}, \texttt{ridge6}) normalized held-out scores through out-of-fold training scores and returned Stop ($-0.0075$); diagnostics showed out-of-fold scores were 2--3$\times$ wider than held-out scores, underweighting MAPA. Switching to label-free \texttt{select\_cdf} still returned Stop for Wave~1 ($-0.0001$) and was fixed for all subsequent screens.
  (2)~\emph{Wave-2 trigger override:} because Wave~1 returned Stop (Ridge6 leave-one-out $-0.00009$), DIVER-1 would not have entered under the original Wave-2 trigger. A descriptive 17-subset Tier~P screen was run anyway, found \texttt{MAPA+DIVER-1} at $+0.0055$ ($+0.0180$ on held-out \texttt{5:0}), and admitted DIVER-1 to E1--E3. The Tier~P gain is therefore optimistically biased as the best of 17 subsets; the one-shot Tier~V veto ($+0.0083$) and one-shot Tier~T test ($+0.0082$) on disjoint evaluation sessions confirm it without selection bias.
\item \textbf{Whole-session preprocessing and contact exclusions.} Both encoders reuse their released session-level preprocessing unchanged: MAPA applies within-shaft mean referencing, whole-session zero-phase 0.5\,Hz FIR high-pass and 60--300\,Hz notch filtering, centered continuous-session STFTs (first Slow frame reaching 250\,ms pre-onset), whole-session median/MAD normalization per contact and frequency bin (spanning the unlabeled test half), and Guard~1 per-session contact exclusions (25 of 1,414 evaluation contacts). DIVER-1 applies whole-session zero-phase 0.5\,Hz high-pass and 60/120/180\,Hz notch filtering and $\pm 200\,\mu\text{V}$ clipping before slicing 0--1\,s windows. Readout standardization uses training rows only.
\item \textbf{Single-seed confirmation.} Tier~V and Tier~T ran DIVER-1 seed~0 only (consistent with single-seed leaderboard practice); Tier~P three-seed spread was $0.6774$--$0.6789$ ($\text{SD}=0.00063$).
\item \textbf{Validation-label use.} Validation labels choose only one ridge penalty per cell and the DIVER early-stopping epoch, as in MAPA and DIVER-1. Half of official validation was used as a one-pass veto set.
\item \textbf{Amendment.} The single-class fallback (Sec.~\ref{sec:amendment}) was added after the first Tier V attempt blocked but before any held-out score was computed; $0/360$ Tier~T cells triggered it.
\item \textbf{Leaderboard export.} After Tier T scoring, the per-session/per-fold test AUROCs required by the leaderboard files were exported by one logged, report-only reread of the Tier T test labels, applied to the already-scored, unmodified predictions (no prediction file is newer than the scoring record). The export reproduced the recorded scores exactly (ensemble $0.7037440086972206$, MAPA $0.695540573759965$; per-session deltas agree to $10^{-12}$). Script: \path{results/submission/export_t_within.py}.
\item \textbf{Pretraining data.} We did no pretraining. MAPA was pretrained on 13 of the 14 permitted sessions (not 5:0) and never on evaluation sessions; DIVER-1 was pretrained without the Brain Treebank. Public PopT and BrainBERT weights (pretrained on 7/12 evaluation sessions) were excluded.
\item \textbf{Prior cross-subject pilots.} A separate cross-subject workstream (\path{07-peer-workstream-audit.md}) ran three negative cross-subject pilots on cross-subject validation labels; none touched within-session data or any test label.
\item \textbf{Scope.} Within-session only; we make no cross-session or cross-subject claim.
\item \textbf{Heterogeneity.} Subject~1 was negative on Tier V; sessions 3:1 and 4:1 were negative on Tier T. Visual tasks (face\_num, frame\_brightness, global\_flow, local\_flow) and volume were negative on Tier T. Gains concentrate in language and onset/delta-volume tasks. No per-task effect survives Holm correction.
\item \textbf{Splits.} Open maintainer PR azaho/neuroprobe\#21 (temporal-block within-session splits) postdates our run; our splits follow \texttt{train\_test\_splits.py} at \texttt{f9b0842} (unchanged at current \texttt{main}, \texttt{b282acd}) and may change under the new splits.
\item \textbf{Effect size.} The gain ($+0.008$ AUROC) is small and rests on 6 subjects.
\end{itemize}

\bibliographystyle{plain}
\bibliography{refs}
\end{document}
